\chapter*{Appendix C: Homotopy Type Theory Foundations of Verified Meaning}
\addcontentsline{toc}{chapter}{Appendix C: Homotopy Type Theory Foundations of Verified Meaning}

\section{Types as Semantic Spaces}

\begin{definition}[Semantic Type]
A semantic type is a homotopy type $A$ whose points represent semantic states and whose paths represent verified equivalences between states.
\end{definition}

\begin{proposition}
The RSVP semantic manifold $M$ is represented by a type $\mathsf{M}$ in which
\[
x =_{\mathsf{M}} y
\]
denotes a homotopy witnessing path-connected semantic equivalence.
\end{proposition}

\section{Identity Types}

\begin{definition}[Identity Type]
For $a,b : A$, the identity type
\[
a =_A b
\]
is the type of equivalences between $a$ and $b$.
\end{definition}

\begin{proposition}
Identity types classify the equivalence of meanings under reconciliation and embedding transformations.
\end{proposition}

\section{Univalence}

\begin{axiom}[Univalence]
For types $A,B$, there is an equivalence
\[
(A \simeq B) \simeq (A =_{\mathsf{Type}} B).
\]
\end{axiom}

\begin{corollary}
Equivalent semantic spaces are identical in all constructions.
\end{corollary}

\section{Higher Inductive Types}

\begin{definition}[Higher Inductive Semantic Object]
A type with:
\begin{itemize}
\item point constructors for primitive semantic atoms,
\item path constructors for rewrite equivalences,
\item higher path constructors for coherence of EHR rules.
\end{itemize}
\end{definition}

\begin{proposition}
A Polyxan hypergraph is represented by a higher inductive type $\mathsf{G}$ with vertices, edges, incidence paths, and coherence cells.
\end{proposition}

\section{Dependent Types and Embeddings}

\begin{definition}[Embedding Family]
A dependent family
\[
\iota : \mathsf{G} \to \mathsf{M}
\]
assigns to each hypergraph cell its embedding in the semantic manifold.
\end{definition}

\begin{proposition}
Curvature-matching and entropy-alignment correspond to dependent path equalities in $\iota$.
\end{proposition}

\section{Modalities}

\begin{definition}[Modality]
A modality is an idempotent monad on $\mathsf{Type}$:
\[
\Box : \mathsf{Type}\to\mathsf{Type},
\quad \eta : A \to \Box A,
\quad \mu : \Box\Box A \to \Box A.
\]
\end{definition}

\begin{proposition}
The operators $\Box_{\mathrm{RSVP}}$ and $\Box_{\mathrm{Polyx}}$ are modalities induced by curvature contraction and EHR normalization.
\end{proposition}

\section{Truncation}

\begin{definition}[n-Truncation]
The truncation $\lVert A\rVert_n$ collapses all paths above level $n$.
\end{definition}

\begin{proposition}
Semantic galaxies correspond to $0$-truncated normal forms inside the homotopy type of all embeddings.
\end{proposition}

\section{H-Propositions and Verified Statements}

\begin{definition}[H-Prop]
A type $P$ is an h-prop if for all $p,q:P$ we have $p=q$.
\end{definition}

\begin{proposition}
Verified modal statements are h-props under $\Box_{\mathrm{RSVP}}$ and $\Box_{\mathrm{Polyx}}$.
\end{proposition}

\section{Pushouts and Reconciliation}

\begin{definition}[Homotopy Pushout]
A pushout of
\[
A \xleftarrow{f} C \xrightarrow{g} B
\]
is a HIT with constructors for $A$, $B$, and paths identifying $f(c)$ with $g(c)$.
\end{definition}

\begin{proposition}
Reconciliation of semantic structures is a homotopy pushout along shared subtypes.
\end{proposition}

\section{Colimits of Rewrite Systems}

\begin{definition}[Rewrite Colimit]
Given a rewrite diagram $\mathsf{G}_0 \to \mathsf{G}_1 \to \ldots$ the colimit $\mathsf{G}_\infty$ is the HIT containing:
\begin{itemize}
\item point constructors for each $\mathsf{G}_i$,
\item path constructors equating successive rewrites.
\end{itemize}
\end{definition}

\begin{proposition}
EHR normal forms correspond to colimits of rewrite diagrams.
\end{proposition}

\section{Synthetic Duality in HoTT}

\begin{definition}[Duality Map]
A duality map is a pair of functions
\[
\mathbb{D} : \mathsf{M} \to \mathsf{G}, \quad
\mathbb{D}^{-1} : \mathsf{G} \to \mathsf{M}
\]
with homotopies expressing inverse equivalence.
\end{definition}

\begin{proposition}
Synthetic duality (Chapter 62) is a HoTT equivalence.
\end{proposition}

\section{Harmonicity as a Path Space Property}

\begin{definition}[Harmonic Embedding]
An embedding $\iota : \mathsf{G} \to \mathsf{M}$ is harmonic if the second derivative vanishes in all infinitesimal directions $D$.
\end{definition}

\begin{proposition}
Harmonic embeddings correspond to fixed points of the $\Box_{\mathrm{RSVP}}$ modality.
\end{proposition}

\section{Quine Stability as Higher Idempotence}

\begin{definition}[Quine-Stable Type]
A type $A$ is quine-stable if the universal quine monad $\mathbb{Q}$ induces an idempotent endomap:
\[
\mathbb{Q}(A) \simeq A.
\]
\end{definition}

\begin{proposition}
The universal media quine $\mathfrak{U}$ is a quine-stable type with contractible higher path structure.
\end{proposition}

\section{Semantic Galaxies as Higher Normal Forms}

\begin{definition}
A semantic galaxy is a 0-truncated, quine-stable, $\Box_{\mathrm{RSVP}}$–$\Box_{\mathrm{Polyx}}$ fixed point.
\end{definition}

\begin{proposition}
Every curvature-reducing evolution yields a semantic galaxy as its homotopy normal form.
\end{proposition}

\section{Collapse Theorem}

\begin{theorem}[Homotopy Collapse]
Let $\mathsf{X}$ be the type of all coupled configurations.  
Repeated application of the modalities yields:
\[
\Box_{\mathrm{RSVP}}\Box_{\mathrm{Polyx}}(\mathsf{X})
  \simeq \mathfrak{U}.
\]
\end{theorem}

\begin{corollary}
The universal media quine is the terminal fixed point of all homotopy evolutions.
\end{corollary}

\section{Summary}

Homotopy type theory provides:
\begin{itemize}
\item identity types for semantic equivalence,
\item HITs for hypergraph structure,
\item modalities for verified invariance,
\item univalence for equivalence of semantic spaces,
\item and homotopy colimits for reconciliation.
\end{itemize}

It forms the proof-relevant foundation of verified meaning.

